
\begin{abstract}

On-policy self-distillation (OPSD) is a promising approach to improve reasoning language
models, but it remains brittle in practice: making it work reliably often requires substantial
engineering effort. We identify a structural source of this difficulty: vanilla OPSD is
precisely the $\beta=1$ member of a broader policy-optimization family, where $\beta$ weights
the KL penalty anchoring the student to a reference policy. This equivalence turns $\beta$ from an implicit value
fixed at one into a controllable regularization parameter, yielding a
more general formulation that trades off proximity to a reference policy against privileged
teacher guidance. We introduce \oursbold and derive its optimal policy as a geometric
interpolation between the reference policy and the privileged teacher. Directly optimizing
this objective with reinforcement learning, however, would be costly and high-variance.
Rather than optimize the RL objective directly, we turn its closed-form solution into a
distillation target. Each value of $\beta$ selects a target along the reference-to-teacher
path, which we implement efficiently by mixing their token-level logits. In this way,
inexpensive distillation approximates the solution of
expensive policy optimization. Return-to-go credit assignment further aligns token updates
with the sequence-level objective while retaining the simplicity of OPSD. Experiments on
mathematical reasoning benchmarks show that \ours{} consistently
outperforms vanilla OPSD, improving optimization stability and downstream reasoning
performance. Our results provide a principled route from self-distillation to policy
optimization and back---without sacrificing the efficiency that makes OPSD practical.

\end{abstract}
